% v11 addition after quintic_prime_support.tex.  No existing source is edited.
% This is an unbounded-prime proof; the companion program only checks examples
% and exhausts the small F_5 completion problem.

\subsection{七辺 support の排除と密な合同模型の完全記述}
\label{ps11:dense-section}

本節でも条件~\eqref{ps10:row}--\eqref{ps10:clutch} を仮定し，
合同計算には basin 面積の正値性を用いない．
前節の七辺 bound は，五つの cores の局所式を同時に使うと強化できる．

\begin{theorem}[密な素数 support]\label{ps11:dense}
任意の奇素数 $\ell$ に対し，$|E(G_\ell)|\ge7$ なら
\[
 \ell=5,\qquad G_\ell\cong K_{2,2,1}.
\]
特に，七辺ちょうどの $G_\ell$ は存在せず，
\begin{equation}\label{ps11:bounds}
 |E(G_\ell)|\le
 \begin{cases}
 2&\ell=3,\\
 8&\ell=5,\\
 6&\ell\ne3,5
 \end{cases}
\end{equation}
となる．$\ell=5$ でも，許される edge 数は $0,\ldots,6,8$ に限られる．

固定された標識付き $K_{2,2,1}$ に対し，対応する $\mathbf F_5$ 上の
row sums，determinants，零分母辺の clutching の解は，後述する
$s=\pm1$ の二つだけである．これらは normal-action 係数の対向相殺も満たす．
これは合同式の解の分類であり，正整数への lifting や大域作用 atlas の実現を含まない．
\end{theorem}

\begin{lemma}[二部 graph に一本を加えた場合]\label{ps11:seven-bipartite}
$G_\ell$ が，parts $\{a,b\},\{c,d,e\}$ の $K_{2,3}$ に
辺 $cd$ を加えた七辺 graph $H$ を含むとする．
このとき $\ell=5$ であり，$ce,de$ のちょうど一方も $G_5$ の辺である．
従って $G_5$ は $K_{2,2,1}$ であり，$H$ そのものを exact support とすることはできない．
\end{lemma}

\begin{proof}
三角形 $acd,bcd$ の符号を $s,t\in\{\pm1\}$ とする．
core $c$ の row $d$ は，最初の二 entries が $s,t$ であり，
第三 entry $-s-t$ も非零である．従って $t=s$ である．
この合同計算の間だけ，全 $p,q$ を同じ $s$ で割る．
$s^2=1$ なので row sums，determinants，零分母辺の clutching は保たれ，
二つの三角形の符号は一に正規化される．

\[
 u=p^c_{ab},\qquad v=p^d_{ab},\qquad
 x=q_{ab},\quad y=q_{ce},\quad z=q_{de}
\]
と置く．core $c$ の rows $a,b,d$ は，列順を $a,b,d,e$ として
\[
 \begin{pmatrix}
 0&u&-1&1-u\\
 -1/u&0&-1&1+1/u\\
 1&1&0&-2
 \end{pmatrix},
\]
core $d$ の対応する rows は，列順 $a,b,c,e$ で
\[
 \begin{pmatrix}
 0&v&1&-1-v\\
 -1/v&0&1&-1+1/v\\
 -1&-1&0&2
 \end{pmatrix}
\]
となる．高分母の局所 branch に接する entries は全て非零なので，
$u,v,1-u,1+u,1-v,1+v$ は非零である．

辺 $ac,ad,bc,bd$ の clutching と各 core の determinants を使うと，
core $a$ の row $b$ と core $b$ の row $a$ はそれぞれ
\begin{align}
 -x&=\frac1u+\frac1v+
       \frac{(1-u)(1+v)}{u+v},\label{ps11:x-a}\\
 -x&=-u-v+\frac{(1+u)(1-v)}{u+v}.
       \label{ps11:x-b}
\end{align}
ここで $u+v\ne0$ は，core $a$ の row $e$ の entry
\[
 p^a_{eb}=-\frac{u+v}{(1-u)(1+v)}
\]
が非零であることから従う．辺 $ae,be$ の clutching により，
core $e$ の rows $a,b$ は列順 $a,b,c,d$ で
\begin{equation}\label{ps11:e-high-rows}
 \begin{pmatrix}
 0&\dfrac{u+v}{(1-u)(1+v)}&-\dfrac1{1-u}&\dfrac1{1+v}\\[6pt]
 \dfrac{u+v}{(1+u)(1-v)}&0&-\dfrac u{1+u}&-\dfrac v{1-v}
 \end{pmatrix}
\end{equation}
となる．その pair $ab$ の determinant は
\[
 (u+v)^2+(1-u^2)(1-v^2)=(uv+1)^2=0
\]
を要求する．体上の計算なので $uv=-1$ である．
\eqref{ps11:x-a} から \eqref{ps11:x-b} を引くと
\[
 0=\frac{u+v}{uv}+(u+v)+\frac{2(v-u)}{u+v}
   =\frac{2(v-u)}{u+v}.
\]
$\ell$ は奇数であるため，$v=u$，従って $u^2=-1$ を得る．
core $c,d$ の残る row sums と \eqref{ps11:x-a} は
\begin{equation}\label{ps11:xyz}
 x=3u,\qquad y=u+\frac12,\qquad z=u-\frac12
\end{equation}
を与える．具体的に，$y,z$ を決める式は
\[
 y=\frac1{1-u}+\frac u{1+u}-\frac12,
 \qquad
 z=-\frac1{1+v}+\frac v{1-v}+\frac12
\]
である．

\eqref{ps11:e-high-rows} の reverse entries を determinants で求め，
残る二つの row sums を使うと，core $e$ の rows $c,d$ は
\[
 \begin{pmatrix}
 1-u&1-u&0&u-5/2\\
 -1-u&-1-u&u+5/2&0
 \end{pmatrix}
\]
となる．最後の pair $cd$ の determinant の左右の差は
\begin{equation}\label{ps11:five-residual}
 \left(u-\frac52\right)\left(u+\frac52\right)
 -\left[\left(u+\frac12\right)\left(u-\frac12\right)-1\right]
 =-5.
\end{equation}
従って $\ell=5$ が必要である．$\mathbf F_5$ 内の $u^2=-1$ の解は
$u=\pm2$ なので，\eqref{ps11:xyz} から $y,z$ のちょうど一方は零，
他方は非零となる．また $x=3u\ne0$ である．
元の符号 $s$ を戻しても零か否かは変わらず，結論を得る．
特に，ここで禁止しているのは exact 七辺 support であり，
$H$ を部分 graph として含む $K_{2,2,1}$ を排除していない．
\end{proof}

次の小さな局所計算を，もう一つの七辺 graph と八辺模型の記述に使う．

\begin{lemma}[三つの零分母を持つ局所合同模型]\label{ps11:three-zero}
$\mathbf F_5$ 上の局所四 branch を $o,b,c,d$ とし，
分母が $(0,r,0,0)$，$r\ne0$ であるとする．
row $o$ が $(p_{ob},p_{oc},p_{od})=(2s,-s,-s)$，$s=\pm1$ なら，
row sums と determinants は
\begin{equation}\label{ps11:outer-matrix}
 r=-s,\qquad
 P=s\begin{pmatrix}
 0&2&-1&-1\\
 2&0&2&2\\
 1&2&0&2\\
 1&2&2&0
 \end{pmatrix}
\end{equation}
を強制する．
\end{lemma}

\begin{proof}
$p_{cd}=st$ と置く．determinants と rows $c,d$ から
\[
 p_{cb}=-s(1+t),\quad p_{db}=s(-1+t^{-1}),\quad
 p_{bc}=\frac{s}{1+t},\quad p_{bd}=\frac{s}{1-t^{-1}}.
\]
全 entries は零分母 branch に接するので非零であり，$t\ne0,\pm1$ である．
従って $t=\pm2$ である．row $b$ の和は，$t=2$ なら $s$，
$t=-2$ なら零である．$r\ne0$ より前者に限られ，表示した値が全て従う．
\end{proof}

\begin{lemma}[$K_1\vee P_4$ の completion]\label{ps11:path-completion}
$G_\ell$ が，vertex $o$ と path $a-b-c-d$ の join を含むなら，
$\ell=5$ であり，さらに辺 $ad$ を含む．従ってその support は
parts $\{a,c\},\{b,d\},\{o\}$ の $K_{2,2,1}$ に等しい．
\end{lemma}

\begin{proof}
定理~\ref{ps10:forbidden} により $\ell=5$ であり，$K_4$ の禁止により
$q_{ac},q_{bd}$ は非零である．同定理の証明で得た core $o$ の行列を
$\mathbf F_5$ で書くと，列・行の順 $a,b,c,d$ に対し
\[
 P^o=s\begin{pmatrix}
 0&1&3&1\\
 4&0&4&2\\
 3&1&0&1\\
 4&2&4&0
 \end{pmatrix},\qquad s\in\{\pm1\}.
\]
core $c$ と core $b$ に補題~\ref{ps11:three-zero} を適用すると，それぞれ
\[
 q_{ac}=-s,\qquad q_{bd}=s
\]
となる．core $b$ の row $a$ は，列 $o,c,d$ に対して
$(-s,-2s,-2s)$ である．辺 $ab$ の clutching を使うと，
core $a$ の row $b$ は $(s,2s,2s)$ となる．
さらに辺 $ao$ の clutching と determinants により，core $a$ では
\[
 p^a_{co}=2s,\quad p^a_{do}=s,\quad
 p^a_{cb}=p^a_{db}=2s.
\]
$y=q_{ad}$ と置くと，rows $c,d$ より
\[
 p^a_{cd}=2s,\qquad p^a_{dc}=-y-3s.
\]
pair $cd$ の determinant は
\[
 2s(-y-3s)=(-s)y-1
\]
を要求する．$\mathbf F_5$ 内ではこの式は $sy=0$ と同値であるから，$y=0$ である．
これが追加辺 $ad$ を与える．他の二辺を加えると $K_4$ となるので，
support は表示した $K_{2,2,1}$ に等しい．
\end{proof}

\begin{proof}[定理~\ref{ps11:dense} の証明]
八辺以上の場合は系~\ref{ps10:counts-gcd} と命題~\ref{ps10:k221} による．
七辺の場合，complement は三辺 graph である．
孤立 vertices を省くと，それは $K_3,K_{1,3},P_4,P_3\sqcup K_2$ のいずれかである．
最初の二つの complements はそれぞれ $B_3$ と $K_4$ を含むため排除される．
後の二つはそれぞれ $K_1\vee P_4$ と補題~\ref{ps11:seven-bipartite} の $H$ なので，
いずれも補題により八本目の辺を強制する．従って七辺ちょうどの場合はない．
三に関する bound は既存の matching 定理から従う．

残る八辺の場合を具体的に記述する．parts を $\{a,b\},\{c,d\},\{o\}$ と固定する．
命題~\ref{ps10:k221} より，ある $s=\pm1$ に対し
\[
 q_{ab}=-s,\qquad q_{cd}=s,\qquad
 q_{ij}=0\quad(ij\ne ab,cd)
\]
となる．core $o$ の行列は，行・列順 $a,b,c,d$ で
\begin{equation}\label{ps11:central-matrix}
 P^o=s\begin{pmatrix}
 0&-2&1&1\\
 -2&0&1&1\\
 -1&-1&0&2\\
 -1&-1&2&0
 \end{pmatrix}.
\end{equation}
core $a,b$ は，他方を $b'$ と書き，行・列順 $(o,b',c,d)$ で
\eqref{ps11:outer-matrix} の行列である．core $c,d$ は同様に，
行・列順 $(o,d',a,b)$ に対し同じ行列の符号 $s$ を $-s$ に置き換えたものである．
補題~\ref{ps11:three-zero} により，これら以外の局所 completion はない．

逆に，表示した二組のデータを代入すると，全ての row sums，determinants，
零分母辺の clutching が成り立つ．normal-action 係数
\[
 D^i_j=-q_{ij}-\sum_{k\ne i,j}p^i_{kj}
\]
は core $o$ で $s(-1,-1,1,1)$，core $a,b$ の上記順序で
$s(1,0,2,2)$，core $c,d$ の上記順序で $-s(1,0,2,2)$ となる．
従って全十辺で $D^i_j+D^j_i=0$ も満たす．
これは，この八辺型を同じ mod~$5$ 条件だけでさらに排除できないことを示す．
\end{proof}

\begin{corollary}[七分母の共通因子の指数]\label{ps11:dense-gcd}
任意に選んだ七本の異なる branches の分母の最大公約数は五を割る．
特に
\[
 \#\{ij:25\mid q_{ij}\}\le6.
\]
$G_5=K_{2,2,1}$ なら，その八辺の分母の最大公約数は正確に五である．
従って，その八辺のうち少なくとも二本では $v_5(q_{ij})=1$ となる．
\end{corollary}

\begin{proof}
七分母に共通する素因数は，定理~\ref{ps11:dense} により五に限られる．
七分母の全てが二十五で割れると仮定する．その七辺は $G_5=K_{2,2,1}$ の
部分 graph であるから，$K_1\vee P_4$ または補題~\ref{ps11:seven-bipartite} の $H$ となる．
以下では $R=\ZZ/25\ZZ$ 上で計算する．
高分母 branch に接する全 entries は determinant が $-1$ なので units である．
三角形の計算から $x^2=1$ を得るが，この環でもその解は $x=\pm1$ に限られる．
実際，$(x-1)(x+1)=0$ で両因子の差は unit の二であるから，
片方は五で割れず，他方は二十五で割れなければならない．

$K_1\vee P_4$ の場合，定理~\ref{ps10:forbidden} の path 計算がそのまま成り立ち，
\[
 p_{ab}=s,\quad p_{bc}=-s,\quad p_{cd}=s,
 \qquad p_{da}=\frac{2s}{3},\quad
 p_{db}=-\frac{s}{2},\quad p_{dc}=-s
\]
となる．分母二と三はこの環の units であり，row $d$ の和は
$-5s/6\ne0$ in $R$ となる．これは row sum に反する．

$H$ の場合も，補題~\ref{ps11:seven-bipartite} の有理式で現れた分母は
全て units であり，\eqref{ps11:x-a}--\eqref{ps11:e-high-rows} は $R$ 上で有効である．
ただし $r=uv+1$ と置いたとき，この環で $r^2=0$ から $r=0$ とは結論しない．
代わりに，二つの $x$ の式の差から
\[
 2uv(v-u)=-(u+v)^2r
\]
を使う．従って $v-u$ は ideal $(r)$ に属し，
\[
 u^2+1=r-u(v-u),\qquad v^2+1=r+v(v-u)
\]
も $(r)$ に属する．よってこれらの平方は零である．
\[
 A=\frac{2u}{1-u^2},\quad B=\frac{2v}{1-v^2},\quad
 C=u-u^{-1}-A,\quad D=v-v^{-1}-B
\]
と置くと
\[
 C-A=\frac{(u^2+1)^2}{u(u^2-1)}=0,\qquad D-B=0,
 \qquad
 A-B=\frac{2(u-v)r}{(1-u^2)(1-v^2)}=0.
\]
core $c,d$ の残る row sums は $y=A+1/2,z=B-1/2$ を与え，
core $e$ の最後の二 entries は $p^e_{cd}=C-5/2$，$p^e_{dc}=D+5/2$ となる．
その determinant の左右の差は
\[
 CD-AB+\frac52(C-D)+\frac12(A-B)-5=-5\ne0
 \quad\text{in }R.
\]
これも矛盾である．従って七分母の gcd は二十五では割れず，五を割る．

二十五で割れる枝が七本以上あれば，そのうち七本を選んで直前の結論に反する．
最後に $G_5=K_{2,2,1}$ の八分母は全て五で割れ，その gcd は五を割るので，
gcd は正確に五である．二十五で割れるものは高々六本であるため，
残る少なくとも二本の五進付値は一となる．
\end{proof}

\begin{remark}[合同分類と整数接着の区別]\label{ps11:modular-scope}
零でない $q_{ij}$ に対する整数整除
$q_{ij}\mid p^i_{jk}+p^j_{ik}$ は，mod~$5$ に落とすとそれだけでは
追加の制約を課さない．上の二解は，この意味での reduced system の全解である．
素数ごとの residue data が存在しても，正の奇整数 $q$，整数 $p$，
全ての整除条件，正の basin 面積，node 配分，有限 neck の作用接着を
同時に持つとは限らない．この定理は任意分母の大域存在も不存在も結論しない．

標識付き $K_{2,2,1}$ は $5\cdot3=15$ 個であり，各々に二つの符号解がある．
付属の標準 Python による検算は，この三十個の合同データについて全式を確認する．
さらに $K_1\vee P_4$ を含む mod~$5$ データの全百六十通りの
分母・符号 assignments を，numerator の全 residue pairs を含めて列挙し，
八辺 completion の二解のみが残ることを確認する．
一般の素数に対する結論はこの有限計算ではなく，
\eqref{ps11:five-residual} を含む上の恒等式の証明による．
\end{remark}
